\documentclass[a4paper,11pt]{article}
\usepackage[utf8]{inputenc}
\usepackage[T1]{fontenc}
\usepackage[french]{babel}
\usepackage{amsmath, amssymb, amsthm}
\usepackage{graphicx}
\usepackage{tcolorbox}
\usepackage{tikz}
\usetikzlibrary{trees,shapes,arrows,positioning,shadows}
\usepackage{listings}
\usepackage{xcolor}
\usepackage{geometry}
\usepackage{hyperref}
\usepackage{pgfplots}
\pgfplotsset{compat=1.18}

\geometry{hmargin=2.5cm,vmargin=2.5cm}

% Configuration des boîtes
\tcbuselibrary{skins,breakable}

\newtcolorbox{conceptbox}[1][]{
  colback=blue!5!white,
  colframe=blue!75!black,
  title={\textbf{💡 Concept Clé}},
  fonttitle=\bfseries,
  enhanced,
  drop shadow,
  breakable,
  #1
}

\newtcolorbox{mathbox}[1][]{
  colback=green!5!white,
  colframe=green!60!black,
  title={\textbf{📐 Fondement Mathématique}},
  fonttitle=\bfseries,
  enhanced,
  drop shadow,
  breakable,
  #1
}

\newtcolorbox{codebox}[1][]{
  colback=gray!5!white,
  colframe=gray!75!black,
  title={\textbf{💻 Implémentation Python}},
  fonttitle=\bfseries,
  breakable,
  #1
}

\newtcolorbox{warningbox}[1][]{
  colback=red!5!white,
  colframe=red!75!black,
  title={\textbf{⚠️ Attention}},
  fonttitle=\bfseries,
  enhanced,
  #1
}

\newtcolorbox{examplebox}[1][]{
  colback=orange!5!white,
  colframe=orange!75!black,
  title={\textbf{📊 Exemple Pratique}},
  fonttitle=\bfseries,
  enhanced,
  breakable,
  #1
}

% Configuration du code
\definecolor{codegreen}{rgb}{0,0.6,0}
\definecolor{codegray}{rgb}{0.5,0.5,0.5}
\definecolor{codepurple}{rgb}{0.58,0,0.82}
\definecolor{backcolour}{rgb}{0.95,0.95,0.92}

\lstdefinestyle{pythonstyle}{
    language=Python,
    backgroundcolor=\color{backcolour},   
    commentstyle=\color{codegreen},
    keywordstyle=\color{magenta},
    numberstyle=\tiny\color{codegray},
    stringstyle=\color{codepurple},
    basicstyle=\ttfamily\footnotesize,
    breakatwhitespace=false,         
    breaklines=true,                 
    captionpos=b,                    
    keepspaces=true,                 
    numbers=left,                    
    numbersep=5pt,                  
    showspaces=false,                
    showstringspaces=false,
    showtabs=false,                  
    tabsize=2
}
\lstset{style=pythonstyle}

\title{\textbf{Isolation Forest} \\ \large Détection d'Anomalies par Isolation}
\author{Cours de Machine Learning Non Supervisé}
\date{\today}

\begin{document}

\maketitle
\tableofcontents
\newpage

\section{Introduction et Motivation}

\subsection{Contexte}
La détection d'anomalies est cruciale dans de nombreux domaines :
\begin{itemize}
    \item \textbf{Maintenance prédictive} : Détecter les pannes avant qu'elles ne surviennent
    \item \textbf{Cybersécurité} : Identifier les intrusions réseau
    \item \textbf{Finance} : Détecter les fraudes bancaires
    \item \textbf{Santé} : Repérer les pathologies rares
\end{itemize}

\subsection{Limites des Approches Classiques}
Les méthodes traditionnelles (K-Means, GMM) modélisent la \textit{normalité} puis considèrent comme anomalies tout ce qui s'en écarte. Problèmes :
\begin{enumerate}
    \item Nécessitent de définir une distribution normale (hypothèse forte)
    \item Coûteuses en calcul pour des données haute dimension
    \item Sensibles aux paramètres (nombre de clusters, etc.)
\end{enumerate}

\begin{conceptbox}
\textbf{Paradigme d'Isolation Forest :}

Au lieu de modéliser le "normal", on cherche directement à \textbf{isoler} les anomalies. L'idée est que les anomalies sont :
\begin{itemize}
    \item \textbf{Rares} (peu nombreuses)
    \item \textbf{Différentes} (valeurs atypiques)
\end{itemize}
$\Rightarrow$ Elles sont donc \textbf{plus faciles à séparer} du reste des données par des partitions aléatoires.
\end{conceptbox}

\section{Principe de Fonctionnement}

\subsection{L'Arbre d'Isolation (iTree)}

Un \textbf{iTree} est un arbre binaire construit de manière aléatoire :

\begin{enumerate}
    \item \textbf{Sélection aléatoire} d'une caractéristique $q$ parmi les $d$ dimensions
    \item \textbf{Choix aléatoire} d'un seuil de coupure $p$ tel que :
    $$\min(X_q) < p < \max(X_q)$$
    \item \textbf{Partition} des données en deux nœuds enfants selon $X_q < p$ ou $X_q \geq p$
    \item \textbf{Récursion} jusqu'à ce que :
    \begin{itemize}
        \item Le nœud ne contienne qu'un seul point (isolé)
        \item Ou la hauteur maximale soit atteinte
    \end{itemize}
\end{enumerate}

\begin{center}
\begin{tikzpicture}[
    level distance=1.5cm,
    level 1/.style={sibling distance=5cm},
    level 2/.style={sibling distance=2.5cm},
    every node/.style={circle, draw, minimum size=8mm, font=\small},
    edge from parent/.style={draw, -latex}
]
\node[fill=blue!20] {Root}
    child {
        node[fill=green!20] {Dense}
        child {node[fill=green!10] {N}}
        child {node[fill=green!10] {N}}
        edge from parent node[left, draw=none, font=\tiny] {$X_1 > 5$}
    }
    child {
        node[fill=red!20] {Sparse}
        child {node[fill=red!50, double] {\textbf{A}} edge from parent node[left, draw=none, font=\tiny] {$X_2 < 3$}}
        child[missing]
        edge from parent node[right, draw=none, font=\tiny] {$X_1 \leq 5$}
    };
\end{tikzpicture}

\textit{Figure 1 : Exemple d'iTree. L'anomalie (A) est isolée en 2 coupes, les points normaux (N) nécessitent plus de profondeur.}
\end{center}

\subsection{La Forêt d'Isolation}

Une \textbf{Isolation Forest} est un ensemble de $t$ iTrees construits sur des sous-échantillons aléatoires des données.

\begin{examplebox}
\textbf{Algorithme de Construction :}
\begin{enumerate}
    \item Fixer le nombre d'arbres $t$ (ex: 100)
    \item Pour chaque arbre $i = 1, \ldots, t$ :
    \begin{itemize}
        \item Tirer aléatoirement $\psi$ points (ex: 256) du dataset
        \item Construire un iTree sur ce sous-échantillon
    \end{itemize}
    \item Retourner la forêt $\{T_1, T_2, \ldots, T_t\}$
\end{enumerate}
\end{examplebox}

\section{Fondements Mathématiques}

\subsection{Longueur du Chemin}

\begin{mathbox}
\textbf{Définition : Longueur du Chemin}

Pour un point $x$ dans un arbre $T$, la \textbf{longueur du chemin} $h(x)$ est le nombre d'arêtes traversées depuis la racine jusqu'à la feuille contenant $x$.

\textbf{Propriété Clé :}
\begin{itemize}
    \item Anomalie $\Rightarrow$ $h(x)$ \textbf{petit} (isolé rapidement)
    \item Point normal $\Rightarrow$ $h(x)$ \textbf{grand} (nécessite beaucoup de coupes)
\end{itemize}
\end{mathbox}

\subsection{Score d'Anomalie}

La longueur moyenne sur tous les arbres doit être normalisée pour être comparable entre datasets de tailles différentes.

\begin{mathbox}
\textbf{Facteur de Normalisation $c(n)$}

Pour un échantillon de taille $n$, $c(n)$ est la longueur moyenne d'un chemin dans un \textit{Binary Search Tree} (BST) non réussi :
$$c(n) = 2H(n-1) - \frac{2(n-1)}{n}$$
où $H(i) = \ln(i) + \gamma$ est le nombre harmonique ($\gamma \approx 0.5772$ constante d'Euler).

\textbf{Approximation :} $c(n) \approx 2 \ln(n-1)$ pour $n$ grand.
\end{mathbox}

\begin{mathbox}
\textbf{Score d'Anomalie $s(x, n)$}

Le score d'anomalie pour un point $x$ est défini par :
$$s(x, n) = 2^{-\frac{E(h(x))}{c(n)}}$$

où :
\begin{itemize}
    \item $E(h(x)) = \frac{1}{t} \sum_{i=1}^{t} h_i(x)$ est la longueur moyenne du chemin sur les $t$ arbres
    \item $c(n)$ normalise par rapport à la taille de l'échantillon
\end{itemize}

\textbf{Interprétation :}
\begin{align*}
s(x, n) \to 1 &\quad \text{Anomalie certaine} \\
s(x, n) \to 0 &\quad \text{Point normal} \\
s(x, n) \approx 0.5 &\quad \text{Indécis (probablement normal)}
\end{align*}
\end{mathbox}

\subsection{Justification Théorique}

\begin{examplebox}
\textbf{Pourquoi $2^{-E(h(x))/c(n)}$ ?}

\begin{enumerate}
    \item Si $E(h(x)) \approx c(n)$ (longueur typique d'un BST) :
    $$s(x, n) = 2^{-1} = 0.5 \quad \text{(comportement aléatoire)}$$
    
    \item Si $E(h(x)) \ll c(n)$ (chemin très court) :
    $$s(x, n) \to 1 \quad \text{(anomalie)}$$
    
    \item Si $E(h(x)) \gg c(n)$ (chemin très long, impossible en pratique) :
    $$s(x, n) \to 0 \quad \text{(normal)}$$
\end{enumerate}
\end{examplebox}

\section{Hyperparamètres et Tuning}

\subsection{Paramètres Principaux (Scikit-Learn)}

\begin{enumerate}
    \item \texttt{n\_estimators} (défaut: 100)
    \begin{itemize}
        \item Nombre d'arbres dans la forêt
        \item $\uparrow$ Plus stable, mais plus lent
        \item Recommandation : 100-200 suffisent généralement
    \end{itemize}
    
    \item \texttt{max\_samples} (défaut: "auto" = 256)
    \begin{itemize}
        \item Taille du sous-échantillon pour chaque arbre
        \item \textbf{Crucial} : Évite le "swamping" (masquage des anomalies par la masse)
        \item Recommandation : 256 pour $n > 256$, sinon $n$
    \end{itemize}
    
    \item \texttt{contamination} (défaut: "auto")
    \begin{itemize}
        \item Proportion attendue d'anomalies
        \item Détermine le seuil de décision sur le score
        \item Exemple : \texttt{contamination=0.05} $\Rightarrow$ 5\% des points avec les scores les plus élevés seront classés comme anomalies
    \end{itemize}
    
    \item \texttt{max\_features} (défaut: 1.0)
    \begin{itemize}
        \item Fraction de features à considérer pour chaque split
        \item Augmente la diversité si $< 1.0$
    \end{itemize}
\end{enumerate}

\begin{warningbox}
\textbf{Piège du Swamping :}

Si on utilise \texttt{max\_samples = n} (tout le dataset), les anomalies minoritaires peuvent être "noyées" dans la masse de points normaux, rendant leur isolation difficile.

$\Rightarrow$ Toujours utiliser un sous-échantillonnage (256 est un bon compromis).
\end{warningbox}

\section{Méthodes d'Évaluation}

\subsection{Cas Supervisé (Labels Disponibles)}

Lorsque les vraies étiquettes sont connues (comme dans \texttt{ai4i2020.csv}) :

\begin{enumerate}
    \item \textbf{Matrice de Confusion}
    $$\begin{array}{c|cc}
     & \text{Prédit Normal} & \text{Prédit Anomalie} \\
    \hline
    \text{Vrai Normal} & TN & FP \\
    \text{Vrai Anomalie} & FN & TP
    \end{array}$$
    
    \item \textbf{Précision (Precision)}
    $$P = \frac{TP}{TP + FP}$$
    Parmi les alertes levées, combien sont vraies ?
    
    \item \textbf{Rappel (Recall / Sensibilité)}
    $$R = \frac{TP}{TP + FN}$$
    Parmi les vraies anomalies, combien ont été détectées ?
    
    \item \textbf{F1-Score}
    $$F_1 = 2 \cdot \frac{P \cdot R}{P + R}$$
    Moyenne harmonique (équilibre Précision/Rappel)
    
    \item \textbf{ROC AUC}
    Aire sous la courbe ROC (Receiver Operating Characteristic)
    \begin{itemize}
        \item Mesure la capacité du modèle à séparer les classes
        \item Indépendant du seuil choisi
    \end{itemize}
\end{enumerate}

\subsection{Cas Non Supervisé (Pas de Labels)}

\begin{itemize}
    \item Analyse de la distribution des scores
    \item Inspection manuelle des points à score élevé
    \item Validation métier (expertise domaine)
\end{itemize}

\section{Exemple Pratique Complet}

\subsection{Dataset : AI4I 2020 Predictive Maintenance}

\begin{codebox}
\begin{lstlisting}
import pandas as pd
import numpy as np
from sklearn.ensemble import IsolationForest
from sklearn.preprocessing import StandardScaler
from sklearn.metrics import classification_report, confusion_matrix, roc_auc_score
import matplotlib.pyplot as plt
import seaborn as sns

# 1. CHARGEMENT DES DONNÉES
df = pd.read_csv("ai4i2020.csv")

# 2. PRÉPARATION
# Sélection des features numériques
features = ['Air temperature [K]', 'Process temperature [K]', 
            'Rotational speed [rpm]', 'Torque [Nm]', 'Tool wear [min]']
X = df[features]
y = df['Machine failure']  # 0 = Normal, 1 = Panne

# Normalisation (important pour IF)
scaler = StandardScaler()
X_scaled = scaler.fit_transform(X)

# 3. ENTRAÎNEMENT DU MODÈLE
# On estime ~3% de pannes dans le dataset
iso_forest = IsolationForest(
    n_estimators=100,
    max_samples=256,
    contamination=0.03,  # 3% d'anomalies attendues
    random_state=42,
    n_jobs=-1
)

# fit_predict retourne 1 (inlier) ou -1 (outlier)
predictions = iso_forest.fit_predict(X_scaled)

# 4. CONVERSION DES PRÉDICTIONS
# IF : 1 = Normal, -1 = Anomalie
# Nos labels : 0 = Normal, 1 = Panne
y_pred = np.where(predictions == -1, 1, 0)

# 5. ÉVALUATION
print("="*60)
print("MATRICE DE CONFUSION")
print("="*60)
cm = confusion_matrix(y, y_pred)
print(cm)

print("\n" + "="*60)
print("RAPPORT DE CLASSIFICATION")
print("="*60)
print(classification_report(y, y_pred, target_names=['Normal', 'Panne']))

# 6. CALCUL DU ROC AUC
# On utilise les scores (decision_function) pour le AUC
scores = iso_forest.decision_function(X_scaled)
# Attention : scores négatifs = anomalies, on inverse
auc = roc_auc_score(y, -scores)
print(f"\nROC AUC Score : {auc:.4f}")

# 7. VISUALISATION DES SCORES
plt.figure(figsize=(12, 5))

plt.subplot(1, 2, 1)
plt.hist(scores[y == 0], bins=50, alpha=0.7, label='Normal', color='blue')
plt.hist(scores[y == 1], bins=50, alpha=0.7, label='Panne', color='red')
plt.xlabel('Anomaly Score (decision_function)')
plt.ylabel('Fréquence')
plt.title('Distribution des Scores IF')
plt.legend()
plt.grid(True, alpha=0.3)

plt.subplot(1, 2, 2)
sns.heatmap(cm, annot=True, fmt='d', cmap='Blues')
plt.title('Matrice de Confusion')
plt.ylabel('Vrai Label')
plt.xlabel('Prédiction')

plt.tight_layout()
plt.show()
\end{lstlisting}
\end{codebox}

\subsection{Résultats Attendus}

\begin{examplebox}
\textbf{Interprétation Typique :}

\begin{itemize}
    \item \textbf{Recall élevé} (70-80\%) : La plupart des pannes sont détectées
    \item \textbf{Précision faible} (10-20\%) : Beaucoup de fausses alarmes
    \item \textbf{Trade-off} : En maintenance, on préfère inspecter trop souvent que rater une panne critique
\end{itemize}

Pour améliorer la précision :
\begin{enumerate}
    \item Ajuster \texttt{contamination} (réduire si trop de FP)
    \item Feature engineering (ajouter des features discriminantes)
    \item Combiner avec des règles métier
\end{enumerate}
\end{examplebox}

\section{Avantages et Limites}

\subsection{Avantages}

\begin{itemize}
    \item \textbf{Efficacité} : Complexité linéaire $O(t \cdot \psi \cdot \log \psi)$ avec $\psi \ll n$
    \item \textbf{Pas d'hypothèse de distribution} : Fonctionne sur données non-gaussiennes
    \item \textbf{Robuste} : Peu sensible aux paramètres
    \item \textbf{Interprétable} : Le score a une signification claire
\end{itemize}

\subsection{Limites}

\begin{itemize}
    \item \textbf{Anomalies locales} : Moins performant que LOF pour détecter des anomalies dans des clusters denses
    \item \textbf{Haute dimension} : Performance peut se dégrader si $d \gg 100$
    \item \textbf{Contamination} : Nécessite d'estimer la proportion d'anomalies
\end{itemize}

\section{Conclusion}

Isolation Forest est une méthode puissante et efficace pour la détection d'anomalies, particulièrement adaptée aux cas où :
\begin{itemize}
    \item Les anomalies sont \textbf{globalement différentes} (pas seulement localement)
    \item Le dataset est \textbf{volumineux} (scalabilité)
    \item On ne connaît pas la distribution des données normales
\end{itemize}

Pour des anomalies plus subtiles (contextuelles, locales), considérer LOF ou des autoencodeurs.

\end{document}
